\documentclass{xmum_icpc_ccpc_article}

\title{厦门大学马来西亚分校算法竞赛兴趣小组\\模板使用说明}
\author{
    \url{https://github.com/Xiamen-University-Malaysia/XMUM_ICPC_CCPC}\\
    \small 执行组长：\href{mailto:zeyu1024@outlook.com}{崔泽禹}}
\date{\today}

\begin{document}

\XMUMSetLogoPath{../logo}
\XMUMSetCoverTypeCN{模板}
\XMUMSetCoverTypeEN{TEMPLATE}
\XMUMSetCoverMainCN{厦门大学马来西亚分校\\算法竞赛兴趣小组}
\XMUMSetCoverMainEN{Xiamen University Malaysia\\Algorithm Contest Interest Group}
\XMUMSetLogoBlue{xiamen_university_logo_blue.png}
\XMUMSetLogoWhite{xiamen_university_logo_white.png}

\XMUMMakeCoverWithBackground
\XMUMMakeCover

\newpage
\maketitle
\tableofcontents
\thispagestyle{titlepage}

\newpage
\XMUMCenterTextPage{这是一个模板说明}

\pagestyle{fancy}
\pagenumbering{arabic}

\section{编译说明}
\noindent 本模板使用 \LaTeX{} 编写，推荐使用两次 \texttt{XeLaTeX} 编译器进行编译，以确保正确生成目录和引用。编译命令如下：
\begin{verbatim}xelatex xmum_icpc_ccpc_article_example.tex && 
xelatex xmum_icpc_ccpc_article_example.tex
\end{verbatim}
\noindent 或使用集成开发环境（如 \href{https://texstudio.sourceforge.net}{TeXstudio}、\href{https://www.overleaf.com/project}{Overleaf}）进行编译，也可在 VS Code 中使用 \href{https://marketplace.visualstudio.com/items?itemName=James-Yu.latex-workshop}{LaTeX Workshop} 插件进行编译，配置过程请参考 \url{https://zhuanlan.zhihu.com/p/166523064}。
\noindent 模板默认添加缩进，若不需要缩进可在导言区添加 \verb|\setlength\parindent{0pt}| 来取消首行缩进。

\section{可调用的接口}
\noindent 本模板类 \texttt{xmum\_icpc\_ccpc\_article.cls} 已封装封面参数、页面组件与常用排版配置。在此章给出可直接调用的接口，可参考 \texttt{xmum\_icpc\_ccpc\_article\_readme.tex} 进行设置与使用。

\subsection{封面与参数设置}
\noindent 以下命令可在 \verb|\begin{document}| 后、生成封面前进行设置：

\begin{table}[H]
    \centering
    \begin{tabular}{@{}ll@{}}
        \toprule
        命令 & 作用 \\
        \midrule
        \texttt{\string\XMUMSetLogoPath\{path\}} & 设置图像资源目录，默认为上级目录 \texttt{../logo} \\
        \midrule
        \texttt{\string\XMUMSetCoverTypeCN\{text\}} & 设置中文封面类型，如：模板、提案、手册 \\
        \texttt{\string\XMUMSetCoverTypeEN\{text\}} & 设置英文封面类型，如：{\footnotesize TEMPLATE, PROPOSAL, MANUAL / HB} \\
        \midrule
        \texttt{\string\XMUMSetCoverMainCN\{text\}} & 设置中文封面主标题， \\
        & 默认为``厦门大学马来西亚分校算法竞赛兴趣小组'' \\
        \texttt{\string\XMUMSetCoverMainEN\{text\}} & 设置英文封面主标题， \\
        & 默认为{\footnotesize “Xiamen University Malaysia Algorithm Contest Interest Group”} \\
        \midrule
        \texttt{\string\XMUMSetLogoBlue\{file\}} & 设置白底封面使用的深色校徽文件， \\
        & 默认为 \texttt{xiamen\_university\_logo\_blue.png} \\
        \texttt{\string\XMUMSetLogoWhite\{file\}} & 设置深底封面使用的浅色校徽文件， \\
        & 默认为 \texttt{xiamen\_university\_logo\_white.png} \\
        \midrule
        \texttt{\string\XMUMSetBuildingBlue\{file\}} & 设置深色建筑底图文件， \\
        & 默认为 \texttt{xmum\_building\_blue.png} \\
        \texttt{\string\XMUMSetBuildingWhite\{file\}} & 设置白色建筑底图文件， \\
        & 默认为 \texttt{xmum\_building\_white.png} \\
        \texttt{\string\XMUMSetBuildingColor\{file\}} & 设置彩色建筑底图文件， \\
        & 默认为 \texttt{xmum\_building.png} \\
        \bottomrule
    \end{tabular}
    \caption{封面参数命令一览}
\end{table}

\noindent 例如：
\begin{verbatim}
\XMUMSetLogoPath{../logo}
\XMUMSetCoverTypeCN{模板}
\XMUMSetCoverTypeEN{TEMPLATE}
\XMUMSetLogoWhite{xiamen_university_logo_white.png}
\XMUMSetBuildingColor{xmum_building.png}
\end{verbatim}

\subsection{页面组件命令}
\noindent 类文件提供了若干可复用页面命令：
\begin{table}[H]
    \centering
    \begin{tabular}{@{}ll@{}}
        \toprule
        命令 & 作用 \\
        \midrule
        \texttt{\string\XMUMMakeCover} & 白色背景封面页 \\
        \texttt{\string\XMUMMakeCoverWithBackground} & 深色背景封面页 \\
        \midrule
        \texttt{\string\XMUMLogoPage} & 白底 LOGO 展示页 \\
        \texttt{\string\XMUMLogoPageWithBackground} & 深底 LOGO 展示页 \\
        \midrule
        \texttt{\string\XMUMSchoolBackgroundPageBlue} & 蓝色建筑背景页 \\
        \texttt{\string\XMUMSchoolBackgroundPageWhite} & 白色建筑背景页 \\
        \texttt{\string\XMUMSchoolBackgroundPageColor} & 彩色建筑背景页 \\
        \midrule
        \texttt{\string\XMUMCenterTextPage\{text\}} & 生成一页居中大字页面 \\
        \midrule
        \texttt{\string\XMUMMarkRealLastPage} & 标记正文真实末页（用于页脚``共 N 页''统计） \\
        \bottomrule
    \end{tabular}
    \caption{页面组件命令一览}
\end{table}

\noindent 例如：
\begin{verbatim}
\XMUMCenterTextPage{这是一个模板说明}
\end{verbatim}

\subsection{末页标注说明}
\noindent 本模板使用 \texttt{\string\zlabel} 来标记正文的真实末页，确保页脚中的``共 N 页''统计正确。请在正文最后一页的合适位置（建议标注在末页的顶部）调用该命令：
\begin{verbatim}
\newpage
\XMUMMarkRealLastPage
\end{verbatim}

\subsection{默认版式与环境}
\begin{itemize}
    \item 文档基类：\texttt{ctexart}（\texttt{11pt, a4paper}）。
    \item 页面边距：上下左右均为 \texttt{2.5cm}。
    \item 代码高亮：已预设 \texttt{listings} 的 \texttt{cpp} 高亮风格，可直接使用 \texttt{lstlisting}。
    \item 页眉页脚：默认启用 \texttt{fancyhdr}，包含顶部横幅与底部页码统计。
\end{itemize}

\newpage
\XMUMMarkRealLastPage

\section{代码示例}
\begin{lstlisting}
#include <bits/stdc++.h>
using namespace std;

typedef pair<int, int> pii;

const int N = 2e5 + 5;

vector<pii> G[N];
int dis[N];
priority_queue<pii, vector<pii>, greater<pii>> pq;

void Dijkstra(int s) {
    memset(dis, 0x3f, sizeof(dis));
    dis[s] = 0; pq.push({0, s});
    while (!pq.empty()) {
        auto [mind, u] = pq.top(); pq.pop();
        if (mind > dis[u]) continue;
        for (auto [v, w] : G[u]) {
            if (dis[v] > dis[u] + w) {
                dis[v] = dis[u] + w;
                pq.push({dis[v], v});
            }
        }
    }
}

int main() {
    int n, m, s;
    cin >> n >> m >> s;
    while (m--) {
        int u, v, w;
        cin >> u >> v >> w;
        G[u].push_back({v, w});
    }
    
    Dijkstra(s);
    
    for (int i = 1; i <= n; i++) {
        cout << dis[i] << " ";
    }
    return 0;
}
\end{lstlisting}

\newpage

\section*{附录一：小组联系方式}
\begin{itemize}
    \item 主要邮箱：\url{xmumcpc2025@outlook.com}
    \item 备用邮箱：\url{xmumcpc2025@foxmail.com}
\end{itemize}

\section*{附录二：小组项目地址}
\begin{itemize}
    \item \url{https://github.com/Xiamen-University-Malaysia/XMUM_ICPC_CCPC}
\end{itemize}

\XMUMSchoolBackgroundPageBlue
\XMUMSchoolBackgroundPageWhite
\XMUMSchoolBackgroundPageColor

\XMUMLogoPage
\XMUMLogoPageWithBackground

\end{document}